\section{Исследование существующих аналогов}

Задача восстановления истории построения трёхмерных моделей по их конечной геометрии (обратный инжиниринг конструктивной истории) является одной из ключевых проблем в области систем автоматизированного проектирования (САПР). Несмотря на то, что коммерческие CAD-системы традиционно используют параметрическое представление с историей построения, при обмене данными между различными платформами или при работе с унаследованными моделями эта информация часто теряется. Как отмечается в современной литературе, распространённые форматы обмена, такие как IGES, не поддерживают хранение истории построения, а даже стандарт STEP, хотя и предусматривает такую возможность, полностью не реализован в большинстве коммерческих систем \cite{kim2023reconstruction}. В результате инженеры получают так называемые «немые» (dumb) модели, представленные только в виде B-Rep (Boundary Representation) или триангулированной геометрии, что существенно затрудняет их редактирование, параметризацию и повторное использование.

Актуальность данной проблемы подтверждается как промышленными потребностями, так и активными исследованиями в этой области. По данным обзора Steininger и соавторов \cite{steininger2025enhancing}, применение методов генеративного искусственного интеллекта в САПР выделяется в четыре основных направления: генерация новых дизайнов, реконструкция существующих, поиск аналогов и модификация. Восстановление истории построения относится к области реконструкции, которая в последние годы демонстрирует наиболее интенсивное развитие благодаря успехам глубокого обучения.

\subsection{Промышленные решения для обратного инжиниринга}

На сегодняшний день наиболее распространённым коммерческим решением для восстановления параметрических признаков из импортированной геометрии является модуль FeatureWorks, входящий в состав САПР SOLIDWORKS \cite{solidxperts2025featureworks}. Данный инструмент позволяет распознавать типовые конструктивные элементы (отверстия, выдавливания, скругления, фаски) в импортированных STEP-, IGES- и Parasolid-файлах, после чего пользователь получает частично восстановленное дерево построения с возможностью редактирования параметров. FeatureWorks поддерживает как автоматический режим распознавания, так и интерактивный, где пользователь вручную указывает элементы для распознавания.

Однако, согласно анализу, представленному на официальном блоге SOLIDWORKS, данный инструмент имеет существенные ограничения \cite{solidxperts2025featureworks}. Во-первых, качество распознавания сильно зависит от сложности геометрии: для литейных форм, поверхностей сложной формы и органических объектов точность резко снижается. Во-вторых, FeatureWorks не восстанавливает эскизные ограничения и связи между параметрами, что ограничивает возможности дальнейшей параметризации. В-третьих, модуль не поддерживает работу со сборочными единицами — каждый компонент необходимо обрабатывать отдельно. Эти ограничения демонстрируют, что даже ведущее промышленное решение не обеспечивает полноценного автоматического восстановления истории построения, особенно для сложных моделей.

Следует отметить, что в экосистеме КОМПАС-3D, которая является целевой средой для разрабатываемого модуля «НейроКОМПАС-3D», на данный момент отсутствует встроенный инструмент, аналогичный FeatureWorks. Это создаёт дополнительную потребность в разработке собственного решения, учитывающего специфику форматов и операционной логики КОМПАС-3D.

\subsection{Академические исследования и перспективные подходы}

Академическое сообщество активно исследует проблему восстановления истории построения с использованием методов машинного обучения, и за последние годы был достигнут значительный прогресс. Обзор современных исследований позволяет выделить несколько ключевых направлений.

\subsubsection{Методы на основе традиционных алгоритмов}

Исторически первыми подходами к решению задачи были алгоритмические методы, основанные на декомпозиции геометрии. Среди них можно выделить методы максимальной декомпозиции объёмов (maximal volume decomposition), предложенные Сакураи \cite{sakurai1998maximal}, а также методы альтернативной суммы объёмов (alternating sum of volumes, ASV) и её модификации ASVP, разработанные Кимом и соавторами \cite{kim2000asvp}. Эти подходы позволяют преобразовывать B-Rep модель в представление в виде конструктивной сплошной геометрии (CSG), где форма описывается как последовательность булевых операций над простыми примитивами.

Однако, как отмечается в недавнем обзоре \cite{kim2023reconstruction}, алгоритмические методы имеют фундаментальные ограничения. Во-первых, они требуют значительных вычислительных ресурсов при работе со сложными моделями и могут не завершиться за приемлемое время. Во-вторых, они не способны распознавать параметры примитивов с высокой точностью. В-третьих, такие методы плохо масштабируются на широкий класс операций и не учитывают семантику конструктивных элементов.

\subsubsection{Методы глубокого обучения для распознавания признаков}

С развитием глубокого обучения исследователи обратились к нейросетевым подходам. Первые попытки использовали свёрточные нейронные сети для анализа воксельных представлений и облаков точек. Однако, как показано в работе Cao и соавторов \cite{cao2020brepnet}, такие аппроксимации неизбежно теряют точность геометрического описания и не позволяют в полной мере использовать топологическую информацию, содержащуюся в B-Rep представлении.

Прорывным направлением стало использование графовых нейронных сетей (GNN) для работы с графами смежности граней, непосредственно извлекаемыми из B-Rep моделей. Работа UV-Net \cite{jayaraman2021uvnet} предложила метод, в котором поверхности граней проецируются в UV-пространство для формирования двумерных сеток, обрабатываемых свёрточными сетями, после чего графовая нейронная сеть учитывает связи между соседними гранями. Аналогичный подход реализован в Hierarchical CADNet \cite{zhang2022hierarchical}, где используется иерархическое представление геометрии.

Люксембургская исследовательская группа под руководством Элоны Дюпон \cite{dupont2025design} представила систематическое исследование подходов к восстановлению истории построения. В своей диссертационной работе, защищённой в 2025 году, автор рассматривает несколько последовательных направлений: от восстановления элементов конструктивной истории из существующих B-Rep представлений до использования иерархических архитектур Transformer для прямого предсказания последовательности операций по облаку точек. Особый интерес представляет предложенный автором подход к использованию больших языковых моделей для генерации исполняемого Python-кода, описывающего процесс построения модели, что позволяет обойти ограничения, связанные с отсутствием стандартизированных форматов представления.

\subsubsection{Генерация параметрических моделей из изображений}

Отдельное направление исследований связано с генерацией CAD-моделей непосредственно из изображений. Работа CADCrafter \cite{chen2025cadcrafter}, представленная на конференции CVPR 2025, демонстрирует возможность создания полных CAD-файлов (в формате STEP) по единственному изображению объекта. Авторы используют комбинацию вариационных автоэнкодеров (VAE) и диффузионных моделей на основе архитектуры Transformer для генерации последовательностей CAD-инструкций. Особенностью подхода является введение механизма проверки компилируемости генерируемого кода с использованием DPO (Direct Preference Optimization), что позволяет существенно повысить долю успешно компилируемых и геометрически корректных моделей.

Это исследование демонстрирует принципиальную возможность использования современных архитектур глубокого обучения для решения задачи, сходной с обратным инжинирингом, хотя и в условиях, когда входные данные представлены в более структурированной форме (изображения), а не в виде трёхмерной геометрии.

\subsection{Проблема датасетов для обучения}

Ключевым препятствием для развития методов машинного обучения в области восстановления истории построения является отсутствие достаточных размеченных наборов данных. Как отмечается в обзоре \cite{steininger2025enhancing}, большинство существующих датасетов либо слишком малы, либо не содержат полной информации о последовательности операций построения.

Важным шагом в этом направлении стал выпуск Autodesk Fusion 360 Gallery Dataset \cite{autodesk2020fusion}, содержащего более 8 тысяч моделей с полной историей построения. Этот датасет, впервые представленный в 2020 году, стал основой для многих последующих исследований, включая работы \cite{dupont2025design} и \cite{chen2025cadcrafter}. Однако, как указывается в корейском исследовании по автоматической генерации датасетов \cite{kim2024automatic}, существующих данных всё ещё недостаточно для обучения моделей, способных обобщаться на широкий класс реальных инженерных задач. Авторы предлагают методы автоматической генерации синтетических датасетов на основе случайного комбинирования параметров признаков с наложением ограничений на реалистичность форм.

В рамках настоящей работы также была проделана значительная работа по формированию датасета: из исходного массива в 82 418 моделей после многоступенчатой фильтрации получен рабочий набор из 21 464 моделей, каждая из которых содержит STEP- и STL-представления, а также JSON-файл с подробным описанием всех операций построения. Это позволяет проводить обучение в supervised режиме и может рассматриваться как вклад в развитие открытых данных для данной предметной области.

\subsection{Анализ проблемы и обоснование необходимости разработки}

Проведённый анализ существующих решений и исследований позволяет выделить несколько ключевых выводов:

\begin{enumerate}
    \item \textbf{Отсутствие универсального промышленного решения.} Даже ведущие CAD-системы, такие как SOLIDWORKS с модулем FeatureWorks, предлагают лишь ограниченные возможности по восстановлению истории построения, работающие преимущественно на простых геометриях и не восстанавливающие полную параметрическую связанность модели.
    
    \item \textbf{Необходимость адаптации под конкретную CAD-систему.} Каждая CAD-платформа имеет собственную систему типов операций, параметров и API, что делает невозможным прямое перенесение решений между системами. Для экосистемы КОМПАС-3D специализированное решение отсутствует.
    
    \item \textbf{Перспективность гибридных подходов.} Современные академические исследования \cite{dupont2025design} \cite{chen2025cadcrafter} демонстрируют, что наиболее эффективными являются гибридные архитектуры, сочетающие нейросетевые методы для распознавания геометрических паттернов и детерминированные алгоритмы для точной параметризации. Этот вывод согласуется с архитектурным решением, принятым в настоящем проекте.
    
    \item \textbf{Ключевая роль датасетов.} Успех применения методов машинного обучения напрямую зависит от наличия качественных размеченных данных. Существующие датасеты либо ограничены по объёму, либо не полностью соответствуют специфике конкретной CAD-системы.
    
    \item \textbf{Проблема точности параметров.} Как показано в исследовании Autodesk \cite{autodesk2020fusion}, нейросетевые методы хорошо справляются с определением типа операции и грубой локализацией, но точное восстановление численных параметров остаётся сложной задачей. Это подтверждает правильность выбранного в проекте подхода, при котором нейросеть отвечает за сегментацию, а точное вычисление параметров переносится на сторону CAD-ядра.
\end{enumerate}

На основе проведённого анализа можно заключить, что разработка программного модуля «НейроКОМПАС-3D» является актуальной и востребованной задачей. Предлагаемое решение, основанное на гибридной архитектуре с использованием нейросетевой сегментации геометрии и детерминированной постобработки средствами КОМПАС-3D, позволяет учесть специфику целевой CAD-системы и преодолеть ограничения существующих подходов.